\section{Proposed Methodology}
\label{sec:methodology}

The methodology is structured into four stages: (1) data acquisition and harmonisation, (2) single-layer preprocessing and quality control, (3) multi-omics integration, and (4) downstream analysis and validation. All analyses will be implemented in Python (v3.10+) and R (v4.3+), using established bioinformatics libraries.

\subsection{Stage 1: Data Acquisition and Harmonisation}

Publicly available multi-omics datasets for LBD, DLB, PDD, and closely related conditions (AD and PD without dementia as comparators) will be obtained from the following repositories:

\begin{itemize}
    \item \textbf{AMP-PD (Accelerating Medicines Partnership -- Parkinson's Disease):} Whole-genome sequencing, whole-blood RNA-seq, and clinical data from PPMI, BioFIND, PDBP, and HBS cohorts. Contains subjects with PD, PDD, and healthy controls.
    \item \textbf{UK Biobank:} Genome-wide genotyping arrays, plasma proteomics (Olink Explore), and linked clinical diagnoses for large-scale genetic and proteomic analyses.
    \item \textbf{Gene Expression Omnibus (GEO):} Post-mortem brain transcriptomics datasets in LBD (e.g., GSE series from frontal cortex, substantia nigra, cingulate cortex). DNA methylation array data (EPIC 850K) from brain tissue EWAS cohorts.
    \item \textbf{NIAGADS Genomics Database:} GWAS summary statistics from the LBD Genomics Consortium \cite{chia2025gwas, chia2021lbd}.
    \item \textbf{MetaboLights / GNPS:} Plasma and CSF metabolomics datasets, including the DLB metabolomics cohort described in \cite{frontiers2024metabolomics}.
\end{itemize}

Data harmonisation will address phenotype labelling inconsistencies, ancestry stratification, and cohort-specific batch effects. The harmonisation workflow is outlined in Algorithm~\ref{alg:harmonisation}.

\begin{algorithm}
\caption{Multi-Omics Data Harmonisation Workflow}
\label{alg:harmonisation}
\begin{algorithmic}[1]
\State \textbf{Input:} Raw omics matrices $\{X_1, X_2, \ldots, X_K\}$ from $K$ data sources; sample metadata tables.
\For{each source $k = 1$ to $K$}
    \State Apply quality control filters (missing rate $< 0.20$, sample call rate $> 0.95$).
    \State Align sample identifiers across omics layers using shared subject IDs.
    \State Infer and correct for ancestry using the first 10 genomic principal components.
    \State Apply ComBat-seq (RNA-seq) or ComBat (methylation arrays, metabolomics) for batch correction.
    \State Normalise features: TMM for RNA-seq; beta-to-M-value for methylation; log-transform + autoscaling for metabolomics.
\EndFor
\State Retain only samples with data present in $\geq 2$ omics layers.
\State \textbf{Output:} Harmonised, batch-corrected omics matrices with a matched sample manifest.
\end{algorithmic}
\end{algorithm}

\subsection{Stage 2: Single-Layer Analysis (Prior to Integration)}

Before integration, each omics layer will be independently analysed to characterise LBD-associated signals and estimate effect sizes:

\begin{itemize}
    \item \textbf{Genomics:} Polygenic risk score (PRS) computation using GWAS summary statistics \cite{chia2025gwas}. Rare variant burden tests for \textit{GBA1} variants using whole-genome sequencing data.
    \item \textbf{Epigenomics:} Differentially methylated position (DMP) and differentially methylated region (DMR) analysis in brain and blood samples using \texttt{limma} and \texttt{DMRcate} (R packages), replicating and extending \cite{nabais2024methylation}.
    \item \textbf{Transcriptomics:} Differential gene expression (DGE) analysis using \texttt{DESeq2}. Gene set enrichment analysis (GSEA) against MSigDB hallmark and Reactome gene sets.
    \item \textbf{Metabolomics:} Multivariate statistical analysis (PLS-DA, OPLS-DA) to identify discriminant metabolite features. Pathway mapping using the \texttt{MetaboAnalyst} platform \cite{frontiers2024metabolomics}.
    \item \textbf{Immunomics:} Re-analysis of publicly available single-cell data \cite{singlecell2024immune} using \texttt{Seurat} (R) for clustering, cell-type annotation, and differential abundance testing.
\end{itemize}

\subsection{Stage 3: Multi-Omics Integration}

Two complementary integration strategies will be applied to exploit the strengths of different methodological paradigms.

\subsubsection{Strategy A: Unsupervised Factor Analysis (MOFA+)}
Multi-Omics Factor Analysis Plus (MOFA+) \cite{argelaguet2020mofa} will be applied to identify latent factors that explain shared and unique variance across omics layers. Each factor represents a coordinated mode of biological variation; factors will be annotated by their top feature weights and tested for association with clinical variables (DLB/PDD status, disease duration, cognitive scores). This approach handles missing data natively, making it well-suited to the incomplete sample overlap expected across layers.

\subsubsection{Strategy B: Similarity Network Fusion (SNF)}
SNF \cite{wang2014snf} will be applied for patient-level integration, constructing per-layer similarity networks and iteratively fusing them into a consensus patient similarity network. Spectral clustering of the fused network will reveal LBD molecular subtypes. The robustness of identified subtypes will be assessed by stability analysis (bootstrap resampling, $n = 1000$) and correlation with clinical and neuropathological variables.

\subsection{Stage 4: Downstream Analysis and Validation}

\subsubsection{Biomarker Model Development}
A supervised machine learning classifier (L1-penalised logistic regression / LASSO) will be trained on integrated multi-omics features to distinguish: (i) LBD vs. AD, (ii) DLB vs. PDD, and (iii) early vs. late LBD. Performance will be evaluated by cross-validated AUROC (area under the receiver operating characteristic curve). A feature importance analysis will yield a minimal discriminant signature.

\subsubsection{Pathway Convergence and Drug Target Analysis}
Convergent pathway scoring will quantify how many independently significant features from each omics layer map onto a given biological pathway (Reactome, KEGG). Mendelian randomisation using GWAS instruments \cite{chia2025gwas, nalls2019pd} will test putative causal relationships between multi-omics features and LBD risk. Top-ranked targets will be queried against the Open Targets Platform and ChEMBL for existing druggability evidence.
